\documentclass[11pt]{article}
\usepackage[margin=1in]{geometry}
\usepackage{booktabs}
\usepackage{longtable}
\usepackage{multirow}
\usepackage{listings}
\usepackage{xcolor}
\usepackage{hyperref}

\lstset{
  basicstyle=\ttfamily\small,
  keywordstyle=\bfseries,
  commentstyle=\color{gray},
  frame=single,
  breaklines=true
}

\title{Updating \texttt{sstubgen} and \texttt{SchemeModula3Types}\\
  for the mscheme Numeric Tower}
\author{Design Report}
\date{March 2026}

\begin{document}
\maketitle

\begin{quote}
\textbf{Status:} Phase~1 of this proposal has been implemented.
See \texttt{numeric-tower-impl.tex} for implementation details
and \texttt{numeric-tower-phase2.tex} for subsequent phases.
\end{quote}

%======================================================================
\section{Background}

The mscheme Scheme interpreter, embedded in Modula-3 via the
\texttt{modula3scheme} package, recently acquired a full numeric tower.
Previously, all Scheme numbers were represented as \texttt{LONGREAL}
(IEEE~754 double-precision).  Now:

\begin{center}
\begin{tabular}{lll}
\toprule
\textbf{Scheme type} & \textbf{M3 type} & \textbf{Description} \\
\midrule
Fixnum   & \texttt{REF INTEGER} (\texttt{SchemeInt.T}) & Exact integer, cached $[-256..256]$ \\
Bignum   & \texttt{Mpz.T}       & GMP arbitrary-precision integer \\
Inexact  & \texttt{REF LONGREAL} (\texttt{SchemeLongReal.T}) & IEEE~754 float \\
\midrule
\multicolumn{3}{l}{\emph{Extended types (library-provided, not part of core tower):}} \\
BigInt   & \texttt{BigInt.T}     & caltech-parser arbitrary-precision integer \\
Mpfr     & \texttt{Mpfr.T}      & MPFR arbitrary-precision float \\
\bottomrule
\end{tabular}
\end{center}

The auto-stub generator \texttt{sstubgen} and the base-type converter
library \texttt{SchemeModula3Types} were written before this change and
assume all numbers are \texttt{LONGREAL}.  This causes two classes of
bugs:

\begin{enumerate}
\item \textbf{Precision loss.}  Integers converted through
  \texttt{LONGREAL} lose precision beyond 53~bits.  For 64-bit
  \texttt{INTEGER} and \texttt{LONGINT}, this is catastrophic for large
  values.
\item \textbf{Type mismatch.}  Generated converters for opaque types
  (e.g.\ \texttt{Mpz.T}) use \texttt{TYPECASE} to match only the exact
  target type.  A Scheme fixnum (\texttt{REF INTEGER}) or bignum
  (\texttt{Mpz.T}) passed where \texttt{BigInt.T} is expected causes a
  runtime error, because the converter does not know these are
  compatible numeric representations.
\end{enumerate}

%======================================================================
\section{Current State: What's Wrong}

\subsection{ToScheme direction (M3 $\to$ Scheme)}

\begin{center}
\begin{tabular}{lll}
\toprule
\textbf{M3 type} & \textbf{Current code} & \textbf{Bug} \\
\midrule
\texttt{INTEGER}  & \texttt{SchemeLongReal.FromI(i)} & Creates inexact float, not exact integer \\
\texttt{CARDINAL} & \texttt{SchemeLongReal.FromI(c)} & Same \\
\texttt{LONGINT}  & \texttt{FLOAT(x\^{}, LONGREAL)} & Catastrophic: loses bits beyond $2^{53}$ \\
\texttt{LONGREAL} & \texttt{SchemeLongReal.FromLR(l)} & Correct \\
\texttt{REAL}     & \texttt{SchemeLongReal.FromLR(...)} & Correct \\
\texttt{EXTENDED} & \texttt{SchemeLongReal.FromLR(...)} & Correct \\
\bottomrule
\end{tabular}
\end{center}

\subsection{ToModula direction (Scheme $\to$ M3)}

\begin{center}
\begin{tabular}{lll}
\toprule
\textbf{M3 type} & \textbf{Current code} & \textbf{Bug} \\
\midrule
\texttt{INTEGER}  & \texttt{SchemeLongReal.FromO} $\to$ \texttt{ROUND} & Lossy: round-trips through double \\
\texttt{CARDINAL} & Same as \texttt{INTEGER} & Same \\
\texttt{LONGINT}  & Same through \texttt{LONGREAL} & Catastrophic loss \\
\texttt{LONGREAL} & \texttt{SchemeLongReal.FromO(l)} & Correct (already handles all types) \\
\texttt{REAL}     & \texttt{FLOAT(FromO, REAL)} & Correct \\
\texttt{EXTENDED} & \texttt{FLOAT(FromO, EXTENDED)} & Correct \\
\texttt{Mpz.T}   & \texttt{TYPECASE} only & Rejects \texttt{REF INTEGER} fixnums \\
\texttt{BigInt.T} & \texttt{TYPECASE} only & Rejects fixnums and \texttt{Mpz.T} \\
\texttt{Mpfr.T}  & \texttt{TYPECASE} only & Rejects all numeric types \\
\bottomrule
\end{tabular}
\end{center}

The floating-point types (\texttt{LONGREAL}, \texttt{REAL},
\texttt{EXTENDED}) are already correct because
\texttt{SchemeLongReal.FromO} already dispatches on all numeric tower
types internally:
\begin{lstlisting}
TYPECASE x OF
  SchemeInt.T(ri) => RETURN FLOAT(ri^, LONGREAL)
| Mpz.T(m)       => RETURN Mpz.get_d(m)
| T(lr)           => RETURN lr^
| BigInt.T(b)     => RETURN BigInt.ToLongReal(b)
END
\end{lstlisting}

\subsection{sstubgen-generated code for opaque types}

For an opaque type like \texttt{Mpz.T}, sstubgen generates:

\begin{lstlisting}
PROCEDURE ToModula_Mpz_T(x : SchemeObject.T) : Mpz.T
    RAISES { Scheme.E } =
  BEGIN
    IF x # NIL AND ISTYPE(x, REF Mpz.T) THEN
      RETURN NARROW(x,REF Mpz.T)^
    END;
    TYPECASE x OF
      Mpz.T(xx) => RETURN xx
    ELSE
      IF SchemeConvertHooks.AttemptConvertToModula(
           TYPECODE(Mpz.T), x) THEN
        RETURN x
      ELSE
        RAISE Scheme.E("Not of type Mpz.T : " & ...)
      END
    END
  END ToModula_Mpz_T;
\end{lstlisting}

The \texttt{TYPECASE} matches only \texttt{Mpz.T} exactly.  A
\texttt{REF INTEGER} (fixnum) or \texttt{REF LONGREAL} is rejected.
The \texttt{SchemeConvertHooks} fallback is available but no hooks are
registered for numeric conversions.

%======================================================================
\section{Design: Correct Conversions}

Every conversion between a Scheme numeric value and an M3 numeric type
should succeed whenever the value is representable in the target type.
The following tables define the correct behavior.

\subsection{ToScheme (M3 $\to$ Scheme)}

\begin{center}
\begin{tabular}{lll}
\toprule
\textbf{M3 type} & \textbf{Scheme result} & \textbf{Implementation} \\
\midrule
\texttt{INTEGER}  & Exact fixnum or bignum & \texttt{SchemeInt.FromI(x)} \\
\texttt{CARDINAL} & Exact fixnum or bignum & \texttt{SchemeInt.FromI(x)} \\
\texttt{LONGINT}  & Exact fixnum or bignum & See \S\ref{sec:longint} \\
\texttt{REAL}     & Inexact float & \texttt{SchemeLongReal.FromLR(FLOAT(x,LONGREAL))} \\
\texttt{LONGREAL} & Inexact float & \texttt{SchemeLongReal.FromLR(x)} \\
\texttt{EXTENDED} & Inexact float & \texttt{SchemeLongReal.FromLR(FLOAT(x,LONGREAL))} \\
\texttt{Mpz.T}   & Exact fixnum or bignum & \texttt{SchemeInt.MpzToScheme(x)} \\
\texttt{BigInt.T} & Exact fixnum or bignum & Convert to Mpz, then demote \\
\texttt{Mpfr.T}  & Inexact float & \texttt{SchemeLongReal.FromLR(Mpfr.GetLR(x))} \\
\bottomrule
\end{tabular}
\end{center}

\subsubsection{LONGINT}\label{sec:longint}

On 64-bit platforms, \texttt{LONGINT} $=$ \texttt{INTEGER} $= 64$-bit
signed.  On 32-bit platforms, \texttt{LONGINT} is 64-bit while
\texttt{INTEGER} is 32-bit.  The correct conversion is: if the value
fits in \texttt{INTEGER}, use \texttt{SchemeInt.FromI}; otherwise,
create an \texttt{Mpz.T} via \texttt{Mpz.NewInt} (which accepts
\texttt{INTEGER}).  For values outside \texttt{INTEGER} range on
32-bit, we need \texttt{Mpz.set\_si64} or equivalent.

In practice, CM3's current 64-bit targets have
\texttt{LONGINT}~$=$~\texttt{INTEGER}, so this is identity.  The
implementation uses \texttt{VAL(x, INTEGER)} which is a no-op on 64-bit
and a checked runtime error on 32-bit if the value is out of range.

\subsection{ToModula (Scheme $\to$ M3)}

\subsubsection{Built-in integer types}

\begin{center}
\begin{tabular}{ll}
\toprule
\textbf{Source Scheme type} & \textbf{Action for target \texttt{INTEGER}} \\
\midrule
\texttt{REF INTEGER} (fixnum) & Direct: \texttt{ri\^{}} \\
\texttt{Mpz.T} (bignum) & Range check via \texttt{fits\_slong\_p}, then \texttt{get\_si} \\
\texttt{REF LONGREAL} (float) & Check integral, range check, \texttt{ROUND} \\
\texttt{BigInt.T} & Falls through to \texttt{SchemeLongReal.FromO} path \\
\bottomrule
\end{tabular}
\end{center}

Same logic for \texttt{CARDINAL} with unsigned range check
($[0..\mathtt{LAST(CARDINAL)}]$, enforced via \texttt{fits\_slong\_p}
followed by a sign check) and \texttt{LONGINT} with appropriate range.

\subsubsection{Built-in floating types}

\texttt{LONGREAL}, \texttt{REAL}, \texttt{EXTENDED}: no change needed.
\texttt{SchemeLongReal.FromO} already handles all numeric tower types
correctly.  Range checking for \texttt{REAL}/\texttt{EXTENDED} is
implicit in the \texttt{FLOAT} conversion (overflow $\to$ $\pm\infty$).

\subsubsection{Library numeric types}

For \texttt{Mpz.T}, the converter accepts any Scheme numeric value
that is representable as an integer:

\begin{center}
\begin{tabular}{ll}
\toprule
\textbf{Source} & \textbf{$\to$ Mpz.T} \\
\midrule
\texttt{REF INTEGER}   & \texttt{Mpz.NewInt(ri\^{})} \\
\texttt{Mpz.T}         & Identity \\
\texttt{REF LONGREAL}  & Check integral, \texttt{ROUND}, \texttt{Mpz.NewInt} \\
\texttt{BigInt.T}       & Falls through to float path \\
\bottomrule
\end{tabular}
\end{center}

%======================================================================
\section{Architecture}

The conversion system has three layers, and each needs targeted changes:

\subsection{Layer 1: \texttt{SchemeModula3Types} (base type converters)}

\textbf{File:} \texttt{m3-scheme/modula3scheme/src/SchemeModula3Types.m3}

This module provides hand-written converters for built-in M3 types.
These are called directly by sstubgen-generated code when the type is
recognized as a ``base type.''

\textbf{Changes:}
\begin{itemize}
\item Add \texttt{IMPORT SchemeInt, Mpz} to the module.
\item \texttt{ToScheme\_INTEGER}: change from \texttt{SchemeLongReal.FromI}
  to \texttt{SchemeInt.FromI}.
\item \texttt{ToScheme\_CARDINAL}: same.
\item \texttt{ToModula\_INTEGER}: rewrite to dispatch on
  \texttt{SchemeInt.T}, \texttt{Mpz.T}, then fall back to
  \texttt{SchemeLongReal.FromO} for floats.
\item \texttt{ToModula\_CARDINAL}: same with unsigned range check.
\end{itemize}

\texttt{ToModula\_LONGINT} (in \texttt{SchemeModula3TypesCM3.m3})
similarly rewritten to dispatch on exact integers first, with the
\texttt{LONGREAL} path as a fallback.

Implemented \texttt{ToModula\_INTEGER}:
\begin{lstlisting}
PROCEDURE ToModula_INTEGER(c : SchemeObject.T)
    : INTEGER RAISES { Scheme.E } =
  BEGIN
    TYPECASE c OF
      SchemeInt.T(ri) => RETURN ri^
    | Mpz.T(m) =>
      IF Mpz.fits_slong_p(m) # 0 THEN
        RETURN Mpz.get_si(m)
      ELSE
        RAISE Scheme.E("INTEGER out of range : "
                        & Stringify(c))
      END
    ELSE
      WITH flt = SchemeLongReal.FromO(c) DO
        IF flt < FLOAT(FIRST(INTEGER), LONGREAL) OR
           flt > FLOAT(LAST(INTEGER), LONGREAL) THEN
          RAISE Scheme.E("INTEGER out of range : "
                          & Stringify(c))
        END;
        WITH int = ROUND(flt) DO
          IF FLOAT(int, LONGREAL) # flt THEN
            RAISE Scheme.E("Not an integer : "
                            & Stringify(c))
          END;
          RETURN int
        END
      END
    END
  END ToModula_INTEGER;
\end{lstlisting}

\subsection{Layer 2: \texttt{Mpz.T} as a recognized base type}

Since \texttt{Mpz} is already a transitive dependency of
\texttt{modula3scheme} (via \texttt{mscheme} $\to$ \texttt{mpz}), we
add \texttt{Mpz.T} as a recognized base type.

\textbf{In \texttt{SchemeModula3Types.i3/.m3}:}
\begin{itemize}
\item Add \texttt{ToScheme\_Mpz\_T} and \texttt{ToModula\_Mpz\_T} procedures.
\item Add \texttt{T\_Mpz\_T} to the \texttt{Type} enum, \texttt{Name}
  array, and \texttt{ProcMode} array (mode: \texttt{Concrete}, since
  \texttt{Mpz.T} is already a reference type).
\end{itemize}

\textbf{In \texttt{TypeTranslator.m3} (sstubgen):}
\begin{itemize}
\item Add \texttt{Mpz.T} to \texttt{AddBasetypes} as a
  \texttt{Type.Opaque} with \texttt{revealedSuperType = Type.refany}
  and name \texttt{Qid(intf="Mpz", item="T")}.
\end{itemize}

\textbf{In \texttt{sstubgen.scm}:}
\begin{itemize}
\item Modify \texttt{is-opaque-basetype} to compare by Qid name rather
  than full structural equality.  The previous code used
  \texttt{equal?}\ on the entire type S-expression, which worked for
  built-in opaques (\texttt{TEXT}, \texttt{MUTEX}) because they resolve
  to singleton \texttt{Type.T} objects (same M3 pointer $\to$ same
  memoized S-expression).  For non-built-in opaques like \texttt{Mpz.T},
  the parser creates a fresh \texttt{Type.Opaque}, so the S-expressions
  are structurally equivalent but identity-distinct.  Name-based matching
  is more robust and matches the existing comment (``special handling for
  opaques, by name'').
\end{itemize}

Implemented converters:
\begin{lstlisting}
PROCEDURE ToScheme_Mpz_T(m : Mpz.T) : SchemeObject.T =
  BEGIN RETURN SchemeInt.MpzToScheme(m) END;

PROCEDURE ToModula_Mpz_T(x : SchemeObject.T)
    : Mpz.T RAISES { Scheme.E } =
  BEGIN
    TYPECASE x OF
      SchemeInt.T(ri) => RETURN Mpz.NewInt(ri^)
    | Mpz.T(m) => RETURN m
    ELSE
      WITH flt = SchemeLongReal.FromO(x) DO
        WITH int = ROUND(flt) DO
          IF FLOAT(int, LONGREAL) # flt THEN
            RAISE Scheme.E("Not an integer : "
                            & Stringify(x))
          END;
          RETURN Mpz.NewInt(int)
        END
      END
    END
  END ToModula_Mpz_T;
\end{lstlisting}

\subsection{Layer 3: \texttt{SchemeConvertHooks} for \texttt{BigInt.T} and \texttt{Mpfr.T}}

\texttt{BigInt} lives in \texttt{caltech-parser/cit\_util} and
\texttt{Mpfr} in \texttt{caltech-other/mpfr}.  These are \emph{not}
dependencies of \texttt{modula3scheme}, so they cannot be added as base
types there.  Instead, we use the existing hook mechanism.

\textbf{Approach:} Each package that exposes \texttt{BigInt.T} or
\texttt{Mpfr.T} to Scheme registers conversion hooks at initialization
time.

\textbf{For BigInt.T}, register hooks during \texttt{BigIntSchemeStubs}
initialization (or in the program's startup code):

\begin{lstlisting}
(* REF INTEGER -> BigInt.T *)
SchemeConvertHooks.RegisterToModula(
  TYPECODE(REF INTEGER),
  TYPECODE(BigInt.T),
  NEW(SchemeConvertHooks.T,
      convert := FixnumToBigInt));

(* Mpz.T -> BigInt.T *)
SchemeConvertHooks.RegisterToModula(
  TYPECODE(Mpz.T),
  TYPECODE(BigInt.T),
  NEW(SchemeConvertHooks.T,
      convert := MpzToBigInt));
\end{lstlisting}

Similarly for \texttt{Mpfr.T}, with converters that accept integer and
float Scheme values and produce \texttt{Mpfr.T} at appropriate
precision.

\textbf{Note on the ToScheme direction:} The sstubgen-generated
\texttt{ToScheme} for opaque types simply returns the object
(\texttt{RETURN x}).  This is acceptable for \texttt{BigInt.T} and
\texttt{Mpfr.T} since they are \texttt{<: REFANY = SchemeObject.T}.
However, Scheme code that receives a \texttt{BigInt.T} will find that
\texttt{(number? x)} returns \texttt{\#t} (since
\texttt{SchemeInt.IsNumber} recognizes \texttt{BigInt.T}), but
arithmetic operations may produce unexpected results.  A
\texttt{ToScheme} hook that converts \texttt{BigInt.T} to the native
exact integer representation would be cleaner.

%======================================================================
\section{Range Checking and Special Values}

\subsection{Integer range checking}

The \texttt{ToModula\_INTEGER} and \texttt{ToModula\_CARDINAL}
converters use a three-tier dispatch:

\begin{enumerate}
\item \textbf{Fixnum} (\texttt{SchemeInt.T}): Direct extraction via
  \texttt{ri\^{}}.  For \texttt{CARDINAL}, a sign check is required
  since fixnums can be negative.  No overflow possible since fixnums
  are \texttt{REF INTEGER} and \texttt{INTEGER} $=$ \texttt{CARDINAL}
  in range for non-negative values.

\item \textbf{Bignum} (\texttt{Mpz.T}): Range check via
  \texttt{Mpz.fits\_slong\_p} (for \texttt{INTEGER}) or
  \texttt{Mpz.fits\_slong\_p} plus sign check (for \texttt{CARDINAL}).
  The GMP \texttt{fits\_slong\_p} function returns non-zero iff the
  value fits in a C \texttt{signed long}, which on 64-bit platforms is
  the same range as \texttt{INTEGER}.

  \textbf{Portability note:} On hypothetical 32-bit platforms where
  \texttt{INTEGER} is 32-bit and C \texttt{long} is also 32-bit,
  \texttt{fits\_slong\_p} correctly checks the 32-bit range.  On
  platforms where C \texttt{long} differs from M3 \texttt{INTEGER}
  (none currently exist in CM3), an additional range check would be
  needed after \texttt{get\_si}.

\item \textbf{Float} (\texttt{REF LONGREAL} and others): Falls through
  to \texttt{SchemeLongReal.FromO} which handles all numeric types.
  The resulting \texttt{LONGREAL} is range-checked against
  \texttt{FLOAT(FIRST(INTEGER), LONGREAL)} /
  \texttt{FLOAT(LAST(INTEGER), LONGREAL)}, then checked for integrality
  via \texttt{FLOAT(ROUND(flt), LONGREAL) = flt}.  This correctly
  rejects NaN (which fails all comparisons), $\pm\infty$ (which fails
  the range check), and non-integral values like \texttt{3.5}.
\end{enumerate}

\subsection{Floating-point special values}

The \texttt{SchemeLongReal.FromO} path correctly handles IEEE~754
special values:

\begin{center}
\begin{tabular}{lll}
\toprule
\textbf{Value} & \textbf{Behavior in integer converter} & \textbf{Behavior in float converter} \\
\midrule
NaN            & Rejected (fails range check) & Passed through \\
$+\infty$      & Rejected (fails range check) & Passed through \\
$-\infty$      & Rejected (fails range check) & Passed through \\
$-0.0$         & Accepted as integer 0        & Passed through as $-0.0$ \\
\bottomrule
\end{tabular}
\end{center}

The treatment of $-0.0$ is correct: \texttt{ROUND(-0.0) = 0} and
\texttt{FLOAT(0, LONGREAL) = -0.0} is false in IEEE~754 strict mode,
but CM3's \texttt{FLOAT} produces \texttt{+0.0} for integer~0, so
the integrality check would actually \emph{reject} $-0.0$.  This is
acceptable: $-0.0$ is an inexact float and should not silently convert
to an exact integer.

\subsection{Mpz.T range checking}

The \texttt{ToModula\_Mpz\_T} converter accepts any exact integer
without range issues (Mpz is arbitrary-precision).  For the float
fallback path, the converter calls \texttt{ROUND(flt)} which returns
an \texttt{INTEGER}, then converts to \texttt{Mpz.T} via
\texttt{Mpz.NewInt}.  This means the float-to-Mpz path is limited
to the \texttt{INTEGER} range.  For truly large float values (beyond
$2^{63}$), the \texttt{ROUND} would fail.  This is acceptable because:
\begin{itemize}
\item Floats beyond $2^{53}$ have lost all sub-integer precision
\item The only sensible path for large integers is through exact
  representations (fixnum or bignum), not through floats
\item A future \texttt{Mpz.set\_d} binding could handle the
  float-to-arbitrary-precision case if needed
\end{itemize}

%======================================================================
\section{Future Numeric Tower Extensions}
\label{sec:future}

The current work addresses the core numeric tower (fixnum, bignum,
inexact float) and integrates \texttt{Mpz.T} as a base type.  Three
additional types merit discussion for future integration:
\texttt{Mpfr.T} (arbitrary-precision float), rational numbers, and
complex numbers.

\subsection{Mpfr.T: Arbitrary-precision floats}

\texttt{Mpfr.T} wraps the MPFR library, providing arbitrary-precision
floating-point arithmetic.  Unlike \texttt{BigInt.T} (which is being
superseded by \texttt{Mpz.T}), \texttt{Mpfr.T} fills a genuine gap
in the numeric tower: there is no other representation for
high-precision inexact real numbers.

\subsubsection{Bidirectional conversion}

\texttt{Mpfr.T} should support full bidirectional conversion:

\begin{center}
\begin{tabular}{lll}
\toprule
\textbf{Direction} & \textbf{Source} & \textbf{Target / Action} \\
\midrule
\multirow{3}{*}{$\to$ Mpfr.T} & REF INTEGER (fixnum) & \texttt{Mpfr.SetInt(m, ri\^{})} \\
 & Mpz.T (bignum) & \texttt{Mpfr.SetMpz(m, mpz)} \\
 & REF LONGREAL (float) & \texttt{Mpfr.SetLR(m, lr\^{})} \\
\midrule
\multirow{1}{*}{Mpfr.T $\to$} & (any context) & \texttt{SchemeLongReal.FromLR(Mpfr.GetLR(m))} \\
\bottomrule
\end{tabular}
\end{center}

\subsubsection{Precision questions}

When converting \emph{to} \texttt{Mpfr.T}, what precision should the
result have?  Options:

\begin{enumerate}
\item \textbf{Default precision}: Use MPFR's global default (typically
  53~bits, matching \texttt{LONGREAL}).  Simple but defeats the purpose
  of using Mpfr in the first place.
\item \textbf{Source-appropriate}: Match the precision to the source ---
  53~bits for \texttt{LONGREAL}, exact for integers (enough bits to
  represent the value).  Better, but the caller likely wants a specific
  precision.
\item \textbf{Caller-specified}: Require the caller to provide the
  precision.  This is cleanest but doesn't fit the auto-converter
  framework, which has no way to pass extra parameters.
\end{enumerate}

The pragmatic choice is option~1 (MPFR default precision) for the
auto-converter, with explicit Scheme-callable procedures for
precision-controlled construction.  The auto-converter serves the
common case where \texttt{Mpfr.T} appears in a Modula-3 interface and
the caller passes a Scheme number; full precision control is available
through the Mpfr Scheme bindings.

\subsubsection{Special values}

MPFR supports NaN, $\pm\infty$, and signed zeros, matching IEEE~754
semantics.  The conversion from \texttt{LONGREAL} preserves these:
\texttt{Mpfr.SetLR} handles NaN and infinities correctly.  The
conversion \emph{from} \texttt{Mpfr.T} to \texttt{LONGREAL} via
\texttt{Mpfr.GetLR} also preserves special values, with overflow
(Mpfr value outside \texttt{LONGREAL} range) producing $\pm\infty$.

\subsubsection{Implementation path}

Since \texttt{Mpfr} is not a dependency of \texttt{modula3scheme},
\texttt{Mpfr.T} conversion should use \texttt{SchemeConvertHooks}
(Layer~3), registered by the package that provides \texttt{Mpfr}
Scheme stubs.  Alternatively, if \texttt{Mpfr.T} becomes common
enough, it could be promoted to a base type (like \texttt{Mpz.T}),
but this would add MPFR as a dependency of \texttt{modula3scheme}.

\subsection{Rational numbers}

A Scheme-standard rational number type (\texttt{exact rational},
e.g.\ $3/7$) is not currently part of mscheme.  If added, it would
likely be represented as:

\begin{lstlisting}
(* In a hypothetical SchemeRational module *)
TYPE T = REF RECORD num, den : Mpz.T END;
  (* or wrapping GMP's mpq_t *)
\end{lstlisting}

\subsubsection{Why exact match, not bidirectional}

Rational numbers should \textbf{not} participate in automatic
bidirectional numeric conversion for the stub generator.  The reasons:

\begin{enumerate}
\item \textbf{No natural M3 target type.}  There is no built-in
  Modula-3 rational type.  Converting $3/7$ to \texttt{INTEGER} (by
  truncation or rounding) is lossy and surprising.  Converting to
  \texttt{LONGREAL} ($\approx 0.4286$) is also lossy.

\item \textbf{Ambiguous intent.}  If a Scheme rational $3/7$ is passed
  where \texttt{LONGREAL} is expected, should the converter truncate,
  round, or reject?  The stub generator cannot know the caller's intent.

\item \textbf{Type safety.}  Modula-3 interfaces that accept
  \texttt{LONGREAL} expect IEEE~754 semantics.  Silently converting
  a rational to a float could mask precision issues that the caller
  should handle explicitly.
\end{enumerate}

\subsubsection{Convenience functions}

Instead of auto-conversion, provide explicit M3-side convenience
functions:

\begin{lstlisting}
(* In SchemeRational or similar *)
PROCEDURE ToLongReal(r : T) : LONGREAL;
  (* Exact division: num/den as LONGREAL *)

PROCEDURE ToInteger(r : T) : INTEGER RAISES { Scheme.E };
  (* Error if not an integer rational *)

PROCEDURE FromInteger(i : INTEGER) : T;
  (* Creates i/1 *)

PROCEDURE FromLongReal(lr : LONGREAL) : T;
  (* Best rational approximation *)
\end{lstlisting}

Scheme code can use these explicitly: \texttt{(rational->real r)} or
\texttt{(exact->inexact r)}.  The auto-generated stubs would treat
\texttt{SchemeRational.T} as an opaque type with exact \texttt{TYPECASE}
matching, which is correct: if a procedure expects a rational, pass a
rational.

\subsection{Complex numbers}

A Scheme complex number type (e.g.\ $3+4i$) would likely be:

\begin{lstlisting}
TYPE T = REF RECORD re, im : LONGREAL END;
  (* or REF RECORD re, im : SchemeObject.T END
     for exact complex like 3/7+2i *)
\end{lstlisting}

\subsubsection{Why exact match, not bidirectional}

Like rationals, complex numbers should use exact \texttt{TYPECASE}
matching, not automatic numeric conversion:

\begin{enumerate}
\item \textbf{Dimension mismatch.}  A complex number has two
  components; an \texttt{INTEGER} or \texttt{LONGREAL} has one.
  Auto-converting $3+4i$ to \texttt{LONGREAL} (by taking the real
  part?\ the magnitude?) is ambiguous and error-prone.

\item \textbf{No natural M3 type.}  Modula-3 has no built-in complex
  type.  Any complex representation is library-defined and should be
  matched exactly.

\item \textbf{Real-valued complex.}  The one tempting case is $3+0i$
  (a complex with zero imaginary part).  Even here, auto-conversion
  is dangerous: it changes the type of the value, which may affect
  dispatch in the M3 code.
\end{enumerate}

\subsubsection{Convenience functions}

\begin{lstlisting}
(* In SchemeComplex or similar *)
PROCEDURE RealPart(c : T) : LONGREAL;
PROCEDURE ImagPart(c : T) : LONGREAL;
PROCEDURE FromReals(re, im : LONGREAL) : T;
PROCEDURE Magnitude(c : T) : LONGREAL;
PROCEDURE Phase(c : T) : LONGREAL;

PROCEDURE IsReal(c : T) : BOOLEAN;
  (* TRUE iff im = 0.0 *)
PROCEDURE ToReal(c : T) : LONGREAL RAISES { Scheme.E };
  (* Error if im # 0.0 *)
\end{lstlisting}

\subsubsection{Interaction with the auto-converter}

When complex numbers appear in Scheme and are passed to M3 stubs
expecting \texttt{LONGREAL}, the auto-generated \texttt{TYPECASE}
will fall through to the \texttt{ELSE} branch (since
\texttt{SchemeComplex.T} is not recognized by
\texttt{SchemeLongReal.FromO}).  This will produce a clear error
message.  If desired, a \texttt{SchemeConvertHooks} registration
could be added to convert real-valued complex numbers to
\texttt{LONGREAL}, but this should be opt-in, not default.

\subsection{Summary of future type integration}

\begin{center}
\begin{tabular}{llll}
\toprule
\textbf{Type} & \textbf{Auto-convert?} & \textbf{Mechanism} & \textbf{Rationale} \\
\midrule
\texttt{Mpfr.T}   & Yes, bidirectional & SchemeConvertHooks & Natural float extension \\
Rational           & No, exact match    & TYPECASE only      & Lossy conversion ambiguous \\
Complex            & No, exact match    & TYPECASE only      & Dimension mismatch \\
\bottomrule
\end{tabular}
\end{center}

The key principle: automatic conversion is appropriate when the source
and target types represent the same mathematical domain (integers,
reals) at different precisions.  It is inappropriate when the types
represent different mathematical structures (rationals $\not\subset$
integers; complex $\not\subset$ reals in a type-safe sense).

%======================================================================
\section{Summary of Changes}

\begin{center}
\begin{longtable}{lll}
\toprule
\textbf{File} & \textbf{Package} & \textbf{Change} \\
\midrule
\endfirsthead
\multicolumn{3}{c}{\emph{(continued)}} \\
\toprule
\textbf{File} & \textbf{Package} & \textbf{Change} \\
\midrule
\endhead
\texttt{SchemeModula3Types.i3} & modula3scheme & Add Mpz.T converters, \\
 & & update comments, \\
 & & add T\_Mpz\_T to Type enum \\
\texttt{SchemeModula3Types.m3} & modula3scheme & Rewrite INTEGER/CARDINAL \\
 & & converters; \\
 & & add Mpz.T converters \\
\texttt{SchemeModula3TypesCM3.m3} & modula3scheme & Rewrite LONGINT converters \\
\texttt{m3makefile} & modula3scheme & Add \texttt{import("mpz")} \\
\texttt{TypeTranslator.m3} & sstubgen & Add Mpz.T to basetypes \\
\texttt{sstubgen.scm} & sstubgen & Name-based opaque matching \\
\texttt{SchemeConvertHooks.m3} & mscheme & No changes needed \\
 & & (existing mechanism suffices) \\
\bottomrule
\end{longtable}
\end{center}

%======================================================================
\section{Immediate Priority: Unblock cspc}

The immediate blocker for the cspc BigInt removal project is that
\texttt{cspgrammar} uses \texttt{BigInt.T} in its interfaces
(\texttt{CspExpression.i3}, \texttt{CspAst.i3}, \texttt{CspType.i3},
\texttt{CspInterval.i3}, etc.).  Replacing \texttt{BigInt.T} with
\texttt{Mpz.T} requires that \texttt{ToModula\_Mpz\_T} accept
\texttt{REF INTEGER} fixnums.

The minimal path to unblock cspc is:
\begin{enumerate}
\item Implement Layer~1 (fix INTEGER/CARDINAL) and Layer~2 (add Mpz.T
  as base type) in \texttt{SchemeModula3Types}.
\item Update \texttt{TypeTranslator.m3} so sstubgen recognizes
  \texttt{Mpz.T}.
\item Modify \texttt{is-opaque-basetype} in \texttt{sstubgen.scm} for
  robust name-based matching of non-built-in opaque basetypes.
\item Rebuild sstubgen, then rebuild cspc (which will regenerate stubs
  with the new Mpz.T converter).
\item Replace \texttt{BigInt.T} $\to$ \texttt{Mpz.T} in cspgrammar and
  cspsimlib interfaces.
\end{enumerate}

Layer~3 (BigInt.T/Mpfr.T hooks) is not needed for cspc since we are
\emph{removing} BigInt.T from the interfaces, not keeping it.  It is
needed for any other packages that still expose BigInt.T or Mpfr.T to
Scheme.

%======================================================================
\section{Verification}

\begin{enumerate}
\item Rebuild modula3scheme: \\
  \texttt{cd \~{}/cm3/cm3 \&\& scripts/do-pkg.sh buildship modula3scheme}
\item Rebuild sstubgen: \\
  \texttt{scripts/do-pkg.sh buildship sstubgen}
\item Test basic number round-tripping in mscheme REPL:
\begin{lstlisting}
> (scheme-modula-conversion-mode 'INTEGER)
;; should return 'Concrete
> (scheme-modula-conversion-mode 'Mpz_T)
;; should return 'Concrete
> (exact? 42)  ;; should be #t
\end{lstlisting}
\item Rebuild cspgrammar, cspsimlib, cspc.
\item Run cspc test suite.
\end{enumerate}

\end{document}
